\documentclass[12pt,a4paper]{article}
\usepackage[utf8]{inputenc}
\usepackage[T1]{fontenc}
\usepackage[slovene]{babel}
\usepackage{amsmath,amssymb}
\usepackage{graphicx}
\usepackage{listings}
\usepackage{xcolor}
\usepackage{geometry}
\usepackage{hyperref}
\usepackage{enumitem}
\usepackage{fancyhdr}
\usepackage{tcolorbox}

\geometry{margin=2.5cm}

\definecolor{codegreen}{rgb}{0,0.6,0}
\definecolor{codegray}{rgb}{0.5,0.5,0.5}
\definecolor{codepurple}{rgb}{0.58,0,0.82}
\definecolor{backcolour}{rgb}{0.95,0.95,0.95}

\lstdefinestyle{javastyle}{
    backgroundcolor=\color{backcolour},
    commentstyle=\color{codegreen},
    keywordstyle=\color{blue}\bfseries,
    numberstyle=\tiny\color{codegray},
    stringstyle=\color{codepurple},
    basicstyle=\ttfamily\small,
    breakatwhitespace=false,
    breaklines=true,
    captionpos=b,
    keepspaces=true,
    numbers=left,
    numbersep=5pt,
    showspaces=false,
    showstringspaces=false,
    showtabs=false,
    tabsize=2,
    language=Java
}

\lstset{style=javastyle}

\newtcolorbox{vprasanje}{
    colback=blue!5!white,
    colframe=blue!75!black,
    title=Mo\v{z}no vpra\v{s}anje na zagovoru,
    fonttitle=\bfseries
}

\newtcolorbox{kljucno}{
    colback=red!5!white,
    colframe=red!75!black,
    title=Klju\v{c}ni koncept,
    fonttitle=\bfseries
}

\pagestyle{fancy}
\fancyhf{}
\rhead{Programiranje 3 -- Projekt}
\lhead{89201088}
\rfoot{\thepage}

\title{\textbf{Simulacija naelektrenih delcev} \\ \large Priprava na zagovor projekta -- Programiranje 3}
\author{Gregor Antonaz -- 89201088}
\date{April 2026}

\begin{document}

\maketitle
\tableofcontents
\newpage

%======================================================================
\section{Pregled problema}
%======================================================================

Simuliramo sistem $n$ naelektrenih delcev v 2D ravnini znotraj omejenega pravokotnika. Vsak delec ima:
\begin{itemize}
    \item Pozicijo $(x, y)$
    \item Hitrost $(v_x, v_y)$
    \item Naboj $c$ (pozitiven ali negativen)
    \item Maso $m = 1.0$
\end{itemize}

\subsection{Coulombova sila (modificirana)}

Sila med delcema $i$ in $j$ je:
\[
F = \frac{|c_i \cdot c_j|}{d^2}
\]
kjer je $d$ evklidska razdalja med njima: $d = \sqrt{(x_j - x_i)^2 + (y_j - y_i)^2}$.

\begin{kljucno}
\textbf{POZOR -- ``spin'' problema:} V na\v{s}i simulaciji velja \textbf{obratna logika} kot v pravi fiziki:
\begin{itemize}
    \item Delca z \textbf{enakim} predznakom naboja se \textbf{privla\v{c}ita}
    \item Delca z \textbf{razli\v{c}nim} predznakom naboja se \textbf{odbijata}
\end{itemize}
To je nasprotno od Coulombovega zakona v fiziki, vendar je tako definirano v nalogi (poglavje 11.2.1, Andrews).
\end{kljucno}

\subsection{Mejna sila}

Pravokotnik ima naboj, ki odbija delce. Ko se delec pribli\v{z}a robu (manj kot 5 enot), deluje nanj odbojna sila:
\[
F_{\text{meja}} = \frac{10.0}{\text{dist}^2}
\]
To prepre\v{c}i, da bi delci zapustili simulacijski prostor.

\subsection{Posodobitev stanja (Eulerjev integracija)}

Vsak cikel:
\begin{enumerate}
    \item Izra\v{c}unamo vse sile na vsak delec
    \item Pospe\v{s}ek: $a = F / m$
    \item Nova hitrost: $v' = v + a$
    \item Nova pozicija: $x' = x + v'$
\end{enumerate}

\begin{vprasanje}
Zakaj se najprej izra\v{c}unajo \textbf{vse} sile, \v{s}ele nato posodobijo pozicije?

\textbf{Odgovor:} \v{C}e bi pozicije posodabljali sproti (med izra\v{c}unom sil), bi nekateri delci \v{z}e imeli nove pozicije, drugi pa \v{s}e stare. Rezultat bi bil nedosleden -- odvisen od vrstnega reda obdelave. Z dvofaznim pristopom zagotovimo, da vsi delci ``vidijo'' enako stanje v istem ciklu.
\end{vprasanje}

%======================================================================
\section{Struktura projekta}
%======================================================================

\begin{verbatim}
src/main/java/com/example/chargedparticles/
  SimulationRunner.java          # vstopna tocka (seq + parallel)
                                 # + skupna simulacijska zanka runCycles()
  DistributedRunner.java         # vstopna tocka za MPJ
  model/
    Particle.java                # podatkovni model delca
  simulation/
    ForceKernel.java             # izracun sil in integracija (skupno jedro)
    Simulation.java              # vmesnik
    AbstractSimulation.java      # skupno stanje vseh treh nacinov
    SequentialSimulation.java    # sekvencna impl.
    ParallelSimulation.java      # vzporedna impl.
    DistributedSimulation.java   # porazdeljena impl. (MPJ)
    SimulationFactory.java       # factory vzorec
    SimulationMode.java          # enum: SEQUENTIAL, PARALLEL, DISTRIBUTED
    SimulationParameters.java    # parametri + branje ukazne vrstice
  ui/
    SimulationUI.java            # Swing GUI
\end{verbatim}

\begin{kljucno}
Vsi trije na\v{c}ini uporabljajo \textbf{isto ra\v{c}unsko jedro} (\texttt{ForceKernel}) nad \textbf{isto podatkovno postavitvijo} in \textbf{isto simulacijsko zanko} (\texttt{SimulationRunner.runCycles}). Razlikujejo se samo v tem, \textbf{kdo} izra\v{c}una \textbf{kateri odsek} delcev. Zato izmerjene razlike v \v{c}asu izvirajo iz na\v{c}ina delitve dela in ne iz razlik v izvedbi fizike.
\end{kljucno}

\subsection{Klju\v{c}ni razredi}

\subsubsection{Particle.java}
Preprost POJO z lastnostmi \texttt{x, y, vx, vy, charge, mass}. Ima stati\v{c}no tovarni\v{s}ko metodo \texttt{randomParticle()} ki sprejme \texttt{Random} objekt za deterministi\v{c}no generiranje.

\begin{vprasanje}
Zakaj uporabljamo \texttt{Random} z seed-om 42?

\textbf{Odgovor:} Da zagotovimo \textbf{reproducibilnost}. Vsi trije na\v{c}ini (sekvencni, vzporedni, porazdeljeni) za\v{c}nejo z \textbf{enako} za\v{c}etno razporeditvijo delcev, kar omogo\v{c}a pravilno primerjavo rezultatov in \v{c}asov izvajanja. \v{C}e seed ne bi bil enak, ne bi mogli potrditi, da implementacije dajejo enake rezultate.
\end{vprasanje}

\subsubsection{Simulation.java (vmesnik)}
Definira pogodbo za vse implementacije:
\begin{lstlisting}
public interface Simulation {
    void performOneCycle();
    void syncToParticles(List<Particle> particles);
    String getDescription();
    default void shutdown() { }
}
\end{lstlisting}

\texttt{performOneCycle()} izvede en cikel nad internim stanjem. \texttt{syncToParticles()} to stanje prepi\v{s}e v seznam \texttt{Particle} objektov -- klicemo ga samo pred izrisom ali izpisom, ne v vro\v{c}i zanki.

\begin{vprasanje}
Zakaj \texttt{performOneCycle()} nima parametrov?

\textbf{Odgovor:} Delci in parametri se ne spreminjajo med tekom, zato jih simulacija prejme \textbf{v konstruktorju} in shrani. Tako v vsakem ciklu ni ponovnega prenosa argumentov niti preverjanja, ali je stanje \v{z}e inicializirano.
\end{vprasanje}

\begin{vprasanje}
Zakaj smo uporabili vmesnik?

\textbf{Odgovor:} Vmesnik omogo\v{c}a \textbf{polimorfizem} -- \texttt{SimulationRunner} dela z \texttt{Simulation} referenco in ne ve, katera implementacija te\v{c}e. To omogo\v{c}a enostavno preklapljanje med na\v{c}ini brez spreminjanja klicne kode. Gre za \textbf{Strategy} na\v{c}rtovalski vzorec.
\end{vprasanje}

\subsubsection{AbstractSimulation.java}
Skupna osnova vseh treh na\v{c}inov. Hrani stanje, ki je pri vseh enako:
\begin{lstlisting}
protected final int n;          // stevilo delcev
protected final double[] state; // 6 doublov na delec
protected final double[] forces;// forces[i*2], forces[i*2+1]
protected final SimulationParameters params;
\end{lstlisting}
Ker je stanje skupno, se podrazredi razlikujejo samo po tem, kako razdelijo delo. Tu je tudi \texttt{syncToParticles()}, ki je pri vseh treh na\v{c}inih enak.

\subsubsection{ForceKernel.java}
Stati\v{c}ni razred, ki vsebuje \textbf{vso fiziko} simulacije. Uporabljajo ga vsi trije na\v{c}ini, zato so ena\v{c}be zapisane na enem samem mestu.

\textbf{Podatkovna postavitev.} Delci niso shranjeni kot objekti, ampak v \textbf{ravnem polju} \texttt{double[]}, 6 vrednosti na delec:
\begin{lstlisting}
// [x, y, vx, vy, charge, mass] -> 6 doublov na delec
// delec i se zacne na odmiku i * 6
\end{lstlisting}

\textbf{Dve operaciji:}
\begin{itemize}
    \item \texttt{computeForces(state, n, start, end, forces, params)} -- izra\v{c}una sile za delce v obmo\v{c}ju $[\text{start}, \text{end})$ proti \textbf{vsem} $n$ delcem, vklju\v{c}no z mejnimi silami
    \item \texttt{integrate(state, start, end, forces)} -- posodobi hitrosti in pozicije za isto obmo\v{c}je
\end{itemize}

\begin{kljucno}
Polje \texttt{forces} je \textbf{globalno indeksirano}: rezultat za delec $i$ gre vedno v \texttt{forces[i*2]} in \texttt{forces[i*2+1]}, ne glede na to, kateri izvajalec ga je izra\v{c}unal. Ker so obmo\v{c}ja razli\v{c}nih izvajalcev disjunktna, lahko ve\v{c} niti hkrati pi\v{s}e v isto polje \textbf{brez zaklepanja}.
\end{kljucno}

Logika predznaka pri Coulombovi sili:
\begin{lstlisting}
double sign = (c1 * c2 >= 0) ? 1.0 : -1.0;
fxSum += sign * magnitude * (dx / r);
fySum += sign * magnitude * (dy / r);
\end{lstlisting}

\begin{kljucno}
\texttt{dx = p2.x - p1.x}. \v{C}e je \texttt{sign = 1.0} (enak naboj), je sila v smeri \textbf{proti} p2 = \textbf{privla\v{c}enje}. \v{C}e je \texttt{sign = -1.0} (razli\v{c}en naboj), je sila \textbf{stran} od p2 = \textbf{odbijanje}.
\end{kljucno}

%======================================================================
\section{Sekvencna implementacija}
%======================================================================

\texttt{SequentialSimulation} -- najenostavnej\v{s}a mo\v{z}na implementacija.

\subsection{Algoritem}

\begin{enumerate}
    \item \textbf{Faza 1 -- Izra\v{c}un sil}: Za vsak delec $i$ izra\v{c}unamo skupno silo od vseh ostalih delcev $j$ in mejne sile. Shranimo v polje \texttt{forces[i][2]}.
    \item \textbf{Faza 2 -- Posodobitev}: Za vsak delec posodobimo hitrost in pozicijo.
\end{enumerate}

\subsection{Kompleksnost}

\begin{itemize}
    \item Faza 1: $O(n^2)$ -- vsak delec proti vsem ostalim
    \item Faza 2: $O(n)$ -- linearen prehod
    \item \textbf{Skupaj na cikel}: $O(n^2)$
    \item \textbf{Skupaj za simulacijo}: $O(n^2 \cdot C)$ kjer je $C$ \v{s}tevilo ciklov
\end{itemize}

Celotna sekvencna implementacija je zato samo dva klica jedra nad celotnim obmo\v{c}jem $[0, n)$:

\begin{lstlisting}
public void performOneCycle() {
    ForceKernel.computeForces(state, n, 0, n, forces, params);
    ForceKernel.integrate(state, 0, n, forces);
}
\end{lstlisting}

\begin{vprasanje}
Zakaj je sekvencna razli\v{c}ica tako kratka?

\textbf{Odgovor:} Ker je vsa fizika v \texttt{ForceKernel}, sekvencna razli\v{c}ica pa samo dolo\v{c}i, da eno samo obmo\v{c}je $[0, n)$ obdela ena nit. Vzporedna in porazdeljena razli\v{c}ica kli\v{c}eta \textbf{isti dve metodi}, le z drugimi mejami obmo\v{c}ja. To je tudi razlog, da vse tri dajo popolnoma enak rezultat.
\end{vprasanje}

%======================================================================
\section{Vzporedna implementacija (Parallel)}
%======================================================================

\texttt{ParallelSimulation} -- uporaba \texttt{ExecutorService} za paralelizacijo izra\v{c}una sil.

\subsection{Strategija}

\begin{kljucno}
Paraleliziramo \textbf{Fazo 1} (izra\v{c}un sil). Faza 2 (posodobitev pozicij) ostane sekvencna -- je $O(n)$ in zanemarljivo hitra v primerjavi s Fazo 1 ki je $O(n^2)$.
\end{kljucno}

\subsection{Razdelitev dela}

Delce razdelimo v \textbf{kose} (chunks):
\begin{lstlisting}
int chunkSize = (n + numThreads - 1) / numThreads;
// nit 0: delci [0, chunkSize)
// nit 1: delci [chunkSize, 2*chunkSize)
// ...
\end{lstlisting}

Vsaka nit izra\v{c}una sile za \textbf{svoj kos} delcev, a pri tem prebere \textbf{vse} delce (ker vsak delec vpliva na vsakega drugega).

\subsection{Sinhronizacija}

Naloge (po ena na kos) zgradimo \textbf{enkrat} v konstruktorju, saj se ne spreminjajo. V vsakem ciklu jih samo oddamo:

\begin{lstlisting}
public void performOneCycle() {
    executor.invokeAll(tasks);   // pregrada: caka na vse niti
    ForceKernel.integrate(state, 0, n, forces);
}
\end{lstlisting}

\texttt{invokeAll()} odda vse naloge in se vrne \textbf{\v{s}ele ko so vse kon\v{c}ane} -- deluje torej kot pregrada (barrier) med Fazo 1 in Fazo 2.

\begin{vprasanje}
Zakaj \texttt{invokeAll()} in ne \texttt{submit()} z \texttt{CountDownLatch}?

\textbf{Odgovor:} \texttt{invokeAll()} \v{z}e vsebuje pregrado, zato \v{s}tevca ni treba ro\v{c}no ustvarjati, zmanj\v{s}evati in \v{c}akati nanj. Manj kode pomeni manj mo\v{z}nosti za napako -- npr. pozabljen \texttt{countDown()} v veji z izjemo bi glavno nit blokiral za vedno.
\end{vprasanje}

\begin{vprasanje}
Zakaj ni \textbf{race condition} pri pisanju v \texttt{forces[i]}?

\textbf{Odgovor:} Vsaka nit pi\v{s}e v \textbf{svoj del} polja \texttt{forces} -- nit 0 pi\v{s}e v indekse $[0, \text{chunkSize})$, nit 1 v $[\text{chunkSize}, 2 \cdot \text{chunkSize})$ itd. Ker se indeksi ne prekrivajo, ni potrebe po zaklepanju.
\end{vprasanje}

\begin{vprasanje}
Zakaj Fazo 2 izvajamo sekvencno?

\textbf{Odgovor:} Faza 2 je $O(n)$ -- linearna. Pri $n = 3000$ je to 3000 operacij. Faza 1 je $O(n^2) = 9.000.000$ operacij. Paralelizacija Faze 2 bi prinesla zanemarljiv prihranaek, hkrati pa bi uvedla kompleksnost sinhronizacije.
\end{vprasanje}

\subsection{Samodejno prilagajanje strojni opremi}

\begin{lstlisting}
this.numThreads = Runtime.getRuntime().availableProcessors();
this.executor = Executors.newFixedThreadPool(numThreads);
\end{lstlisting}

\texttt{availableProcessors()} vrne \v{s}tevilo logi\v{c}nih jeder. Na mojem ra\v{c}unalniku je to 10.

\begin{vprasanje}
Zakaj \texttt{FixedThreadPool} in ne \texttt{CachedThreadPool}?

\textbf{Odgovor:} \texttt{FixedThreadPool} ohranja konstantno \v{s}tevilo niti skozi celotno simulacijo. Ne ustvarja novih niti za vsak cikel -- niti se ponovno uporabijo. \texttt{CachedThreadPool} bi lahko ustvaril preve\v{c} niti in povzro\v{c}il overhead.
\end{vprasanje}

\subsection{Pohitritev}

Idealna pohitritev z $T$ nitmi je $T\times$. V praksi je manj\v{s}a zaradi:
\begin{itemize}
    \item \textbf{Overhead} -- oddajanje nalog in pregrada \texttt{invokeAll} v vsakem ciklu
    \item \textbf{Neenakomerna razdelitev} -- \v{c}e $n$ ni deljiv s \v{s}tevilom niti
    \item \textbf{Amdahlov zakon} -- Faza 2 je sekvencna in omejuje maksimalno pohitritev
\end{itemize}

\begin{kljucno}
\textbf{Amdahlov zakon:} \v{C}e je dele\v{z} sekvencne kode $s$, je maksimalna pohitritev:
\[
S = \frac{1}{s + \frac{1-s}{p}}
\]
V na\v{s}em primeru je $s \approx \frac{O(n)}{O(n^2)} \approx \frac{1}{n}$, kar je zelo majhno -- zato se vzporedna verzija dobro skalira.
\end{kljucno}

%======================================================================
\section{Porazdeljena implementacija (MPJ Express)}
%======================================================================

\texttt{DistributedSimulation} + \texttt{DistributedRunner} -- uporaba MPI (Message Passing Interface) preko MPJ Express.

\subsection{Kaj je MPJ Express?}

MPJ Express je Java implementacija MPI standarda. Omogo\v{c}a komunikacijo med ve\v{c} procesi, ki te\v{c}ejo na istem ali razli\v{c}nih ra\v{c}unalnikih.

\begin{vprasanje}
Kaj je razlika med vzporedno in porazdeljeno implementacijo?

\textbf{Odgovor:}
\begin{itemize}
    \item \textbf{Vzporedna}: Ve\v{c} \textbf{niti} v enem procesu, deljeni pomnilnik. Niti direktno berejo isti \texttt{List<Particle>}.
    \item \textbf{Porazdeljena}: Ve\v{c} lo\v{c}enih \textbf{procesov}, vsak s svojim pomnilnikom. Procesi morajo \textbf{eksplicitno komunicirati} (MPI) za izmenjavo podatkov.
\end{itemize}
\end{vprasanje}

\subsection{Zakaj lo\v{c}en DistributedRunner?}

MPJ Express zahteva \texttt{MPI.Init(args)} na samem za\v{c}etku programa. Vsi procesi zazenejo \textbf{isti program} (SPMD -- Single Program, Multiple Data). Zato imamo lo\v{c}en vstopni razred, ki se za\v{z}ene z:
\begin{verbatim}
mpjrun.sh -np 4 com.example.chargedparticles.DistributedRunner
\end{verbatim}

\subsection{Podatkovne strukture za MPI}

MPI ne more po\v{s}iljati Java objektov (\texttt{Particle}). Zato pretvorimo vse delce v \textbf{ravno polje} \texttt{double[]}:

\begin{lstlisting}
// 6 vrednosti na delec: x, y, vx, vy, charge, mass
double[] state = new double[n * 6];

// Delec i je na odmiku i * 6:
// state[i*6 + 0] = x
// state[i*6 + 1] = y
// state[i*6 + 2] = vx
// state[i*6 + 3] = vy
// state[i*6 + 4] = charge
// state[i*6 + 5] = mass
\end{lstlisting}

\begin{vprasanje}
Zakaj ne serializiramo \texttt{Particle} objektov?

\textbf{Odgovor:} Serializacija Java objektov je po\v{c}asna (overhead marshalling/unmarshalling). Ravno \texttt{double[]} polje je najhitrej\v{s}i na\v{c}in za MPI komunikacijo -- MPI.DOUBLE je primitiven tip, ki se kopira direktno brez pretvorb.
\end{vprasanje}

\subsection{Razdelitev delcev med procese}

Delce razdelimo \v{c}im bolj enakomerno:
\begin{lstlisting}
// Zacetni indeks delcev za proces r.
// Prvih (n mod size) procesov dobi en delec vec.
private int startIndexOf(int r) {
    int base = n / size;
    int remainder = n % size;
    return (r < remainder)
        ? r * (base + 1)
        : remainder * (base + 1) + (r - remainder) * base;
}

// Moje obmocje je preprosto:
myStart = startIndexOf(rank);
myEnd   = startIndexOf(rank + 1);
\end{lstlisting}

Primer z 10 delci in 3 procesi:
\begin{center}
\begin{tabular}{|c|c|c|}
\hline
\textbf{Rank} & \textbf{Delci} & \textbf{\v{S}tevilo} \\
\hline
0 & [0, 4) & 4 \\
1 & [4, 7) & 3 \\
2 & [7, 10) & 3 \\
\hline
\end{tabular}
\end{center}

\subsection{Algoritem enega cikla}

\begin{enumerate}
    \item \textbf{Faza 1}: Vsak proces izra\v{c}una sile za \textbf{svoje} delce. Za izra\v{c}un potrebuje podatke \textbf{vseh} delcev (ki jih \v{z}e ima v \texttt{state}).

    \item \textbf{Faza 2}: Vsak proces posodobi pozicije in hitrosti \textbf{samo svojih} delcev v \texttt{state}.

    \item \textbf{Faza 3 -- Allgatherv}: Vsi procesi si izmenjajo posodobljene podatke.
\end{enumerate}

\subsection{Allgatherv -- klju\v{c}na MPI operacija}

\begin{kljucno}
\texttt{Allgatherv} = ``vsi zberejo vse, razli\v{c}ne koli\v{c}ine''.

Vsak proces po\v{s}lje svoj del podatkov in prejme dele od vseh ostalih. Po klicu ima vsak proces \textbf{celoten posodobljen} \texttt{state}.
\end{kljucno}

\begin{lstlisting}
MPI.COMM_WORLD.Allgatherv(
    state, sendOffset, sendCount, MPI.DOUBLE,  // posiljam
    state, 0, recvCounts, displacements, MPI.DOUBLE  // prejemam
);
\end{lstlisting}

Parametri:
\begin{itemize}
    \item \texttt{sendOffset} -- za\v{c}etek mojih podatkov v polju
    \item \texttt{sendCount} -- koliko doublov po\v{s}iljam
    \item \texttt{recvCounts[]} -- koliko doublov pri\v{c}akujem od vsakega procesa
    \item \texttt{displacements[]} -- kje v sprejemnem polju se za\v{c}nejo podatki vsakega procesa
\end{itemize}

\begin{figure}[h]
\centering
\begin{verbatim}
Pred Allgatherv:
  Rank 0: [x0 y0 ... | stari  | stari ]  <- posodobil svoj del
  Rank 1: [stari  | x4 y4 ... | stari ]  <- posodobil svoj del
  Rank 2: [stari  | stari  | x7 y7 ...]  <- posodobil svoj del

Po Allgatherv:
  Rank 0: [x0 y0 ... | x4 y4 ... | x7 y7 ...]  <- vsi novi
  Rank 1: [x0 y0 ... | x4 y4 ... | x7 y7 ...]  <- vsi novi
  Rank 2: [x0 y0 ... | x4 y4 ... | x7 y7 ...]  <- vsi novi
\end{verbatim}
\caption{Vizualizacija Allgatherv operacije}
\end{figure}

\begin{vprasanje}
Zakaj ne uporabimo Scatter + Gather namesto Allgatherv?

\textbf{Odgovor:} Pri Scatter/Gather bi samo rank 0 (master) imel celoten rezultat in bi ga moral ponovno razpo\v{s}ljati (Broadcast). Z \texttt{Allgatherv} \textbf{vsi} procesi hkrati prejmejo celotno posodobljeno stanje -- ena operacija namesto dveh. To je u\v{c}inkovitej\v{s}e.
\end{vprasanje}

\begin{vprasanje}
Zakaj potrebuje vsak proces \textbf{vse} delce, ne le svoje?

\textbf{Odgovor:} Ker sila na delec $i$ je vsota sil od \textbf{vseh} ostalih $n-1$ delcev. Delec na rank 0 \v{c}uti silo od delcev na rank 1, rank 2, itd. Brez poznavanja pozicij vseh delcev izra\v{c}un ni mo\v{z}en. To je inherentna lastnost N-body problema.
\end{vprasanje}

\subsection{Overhead porazdeljene verzije}

Porazdeljena verzija je po\v{c}asnej\v{s}a od vzporedne zaradi:
\begin{enumerate}
    \item \textbf{Allgatherv vsak cikel} -- kopiranje $6n$ doublov med vsemi procesi
    \item \textbf{MPI overhead} -- inicializacija, interna sinhronizacija
    \item \textbf{Lo\v{c}en pomnilnik} -- ni mo\v{z}nosti deljenja podatkov brez komunikacije
\end{enumerate}

Vendar ima porazdeljena verzija prednost: lahko te\v{c}e na \textbf{ve\v{c} ra\v{c}unalnikih}, kar omogo\v{c}a skaliranje prek zmogljivosti enega stroja.

%======================================================================
\section{Primerjava implementacij}
%======================================================================

\begin{center}
\begin{tabular}{|l|c|c|c|}
\hline
& \textbf{Sekvencna} & \textbf{Vzporedna} & \textbf{Porazdeljena} \\
\hline
Enot izvajanja & 1 nit & $T$ niti & $P$ procesov \\
Pomnilnik & deljen & deljen & lo\v{c}en \\
Komunikacija & / & deljen pomnilnik & MPI Allgatherv \\
Sinhronizacija & / & \texttt{invokeAll} pregrada & \texttt{Allgatherv} \\
Vstopna to\v{c}ka & SimulationRunner & SimulationRunner & DistributedRunner \\
Zagon & \texttt{java ...} & \texttt{java ...} & \texttt{mpjrun.sh -np P} \\
\hline
\end{tabular}
\end{center}

\subsection{Meritve}

Meritve zaganjamo z \texttt{./benchmark.sh}. Vsaka nastavitev se izvede ve\v{c}krat, poro\v{c}amo najkraj\v{s}i \v{c}as (druga dejavnost na ra\v{c}unalniku lahko \v{c}as samo pove\v{c}a, nikoli zmanj\v{s}a).

\textbf{2000 delcev, 1000 ciklov:}
\begin{center}
\begin{tabular}{|l|c|c|}
\hline
\textbf{Na\v{c}in} & \textbf{\v{C}as [s]} & \textbf{Pohitritev} \\
\hline
Sekvencna & 6.94 & 1.0$\times$ \\
Porazdeljena (4 proc.) & 2.56 & 2.7$\times$ \\
Porazdeljena (10 proc.) & 2.39 & 2.9$\times$ \\
Vzporedna (10 niti) & \textbf{1.45} & \textbf{4.8$\times$} \\
\hline
\end{tabular}
\end{center}

\textbf{3000 delcev, 1000 ciklov:}
\begin{center}
\begin{tabular}{|l|c|c|}
\hline
\textbf{Na\v{c}in} & \textbf{\v{C}as [s]} & \textbf{Pohitritev} \\
\hline
Sekvencna & 15.67 & 1.0$\times$ \\
Porazdeljena (4 proc.) & 7.72 & 2.0$\times$ \\
Porazdeljena (10 proc.) & 5.42 & 2.9$\times$ \\
Vzporedna (10 niti) & \textbf{4.16} & \textbf{3.8$\times$} \\
\hline
\end{tabular}
\end{center}

\begin{kljucno}
\textbf{Vzporedna je vedno najhitrej\v{s}a na enem ra\v{c}unalniku}, in to je pri\v{c}akovano: obe razli\v{c}ici razdelita enako delo med enako \v{s}tevilo izvajalcev, a porazdeljena mora vsak cikel \v{s}e prekopirati stanje med procesi.

Zanimivej\v{s}i je \textbf{trend}. Zaostanek porazdeljene razli\v{c}ice (10 proc.) za vzporedno pada z velikostjo problema:
\begin{center}
\begin{tabular}{|c|c|}
\hline
$n$ & Porazdeljena po\v{c}asnej\v{s}a za \\
\hline
500  & 278\,\% \\
1000 & 71\,\% \\
2000 & 64\,\% \\
3000 & 30\,\% \\
\hline
\end{tabular}
\end{center}
Razlog: komunikacija raste kot $O(n)$, ra\v{c}unanje pa kot $O(n^2)$. Ve\v{c}ji kot je problem, manj\v{s}i dele\v{z} \v{c}asa gre za komunikacijo.
\end{kljucno}

\begin{vprasanje}
Zakaj je pri $n = 500$ porazdeljena z 10 procesi po\v{c}asnej\v{s}a od \textbf{sekvencne} (pohitritev 0{,}72$\times$)?

\textbf{Odgovor:} Pri 500 delcih je ra\v{c}unanja na cikel tako malo, da 1000 klicev \texttt{Allgatherv} prevlada nad celotnim \v{c}asom izvajanja. Komunikacijski stro\v{s}ek je ve\v{c}ji od prihranka pri ra\v{c}unanju. To je najlep\v{s}i primer v na\v{s}ih podatkih, da paralelizacija ni brezpla\v{c}na.
\end{vprasanje}

\begin{vprasanje}
Zakaj je pri $n = 500$ bolje uporabiti 4 procese kot 10?

\textbf{Odgovor:} Manj procesov pomeni cenej\v{s}o kolektivno operacijo. Pri $n = 500$ so 4 procesi hitrej\v{s}i (0{,}25\,s proti 0{,}61\,s), pri $n = 3000$ pa 10 procesov (5{,}42\,s proti 7{,}72\,s) -- takrat je ra\v{c}unanja dovolj, da se dodatni procesi izpla\v{c}ajo.
\end{vprasanje}

%======================================================================
\section{Na\v{c}rtovalski vzorci}
%======================================================================

\begin{enumerate}
    \item \textbf{Strategy} -- vmesnik \texttt{Simulation} z razli\v{c}nimi implementacijami. \texttt{SimulationRunner} ne ve, katera implementacija te\v{c}e.

    \item \textbf{Factory} -- \texttt{SimulationFactory.createSimulation(mode)} ustvari pravilno implementacijo glede na izbrani na\v{c}in.

    \item \textbf{SPMD} (Single Program, Multiple Data) -- pri porazdeljeni verziji vsi procesi zazenejo isti program (\texttt{DistributedRunner}), a obdelajo razli\v{c}ne podatke glede na svoj rang.
\end{enumerate}

%======================================================================
\section{Zagon programa}
%======================================================================

\subsection{Zagonske skripte}

Vsi trije na\v{c}ini se za\v{z}enejo z eno skripto na na\v{c}in. Skripte same poi\v{s}\v{c}ejo \texttt{JAVA\_HOME} in \texttt{MPJ\_HOME} (prek skupne \texttt{setup\_env.sh}) in prevedejo projekt.

\begin{verbatim}
./run_sequential.sh                  # sekvencno, z graficnim vmesnikom
./run_parallel.sh                    # vzporedno
./run_distributed.sh 4               # porazdeljeno, 4 MPI procesi

# Meritev brez izrisa
./run_distributed.sh 4 --particles 2000 --cycles 1000 --ui false

# Vse meritve iz porocila naenkrat
./benchmark.sh
\end{verbatim}

\subsection{Argumenti}
Opcije so pri vseh treh skriptah \textbf{enake} -- bere jih ista metoda \texttt{SimulationParameters.fromArgs()}.
\begin{center}
\begin{tabular}{|l|l|l|}
\hline
\textbf{Argument} & \textbf{Opis} & \textbf{Privzeto} \\
\hline
\texttt{--mode} & sequential $|$ parallel & sequential \\
\texttt{--particles N} & \v{s}tevilo delcev & 400 \\
\texttt{--cycles N} & \v{s}tevilo ciklov & 1000 \\
\texttt{--ui true|false} & graficni vmesnik & true \\
\texttt{--window W H} & velikost okna & 800 600 \\
\texttt{--bounds x1 x2 y1 y2} & meje prostora & 0 800 0 600 \\
\texttt{--fps N} & hitrost osve\v{z}evanja & 60 \\
\hline
\end{tabular}
\end{center}

\subsection{Preverjanje pravilnosti}
Ker vsi na\v{c}ini startajo iz istega seeda, morajo dati \textbf{enak} rezultat:
\begin{verbatim}
./run_sequential.sh    --particles 500 --cycles 300 --ui false
./run_parallel.sh      --particles 500 --cycles 300 --ui false
./run_distributed.sh 4 --particles 500 --cycles 300 --ui false
\end{verbatim}
Izpisanih pet delcev na koncu se mora v vseh treh primerih ujemati do zadnje decimalke.

\subsection{Porazdeljeni na\v{c}in iz graficnega vmesnika}
V oknu izberemo na\v{c}in \textbf{Porazdeljena (MPI)}; pojavi se polje \textbf{Procesi}. Klik na \textbf{Start} pока\v{z}e ukaz, ki bo izveden, in po potrditvi zaz\v{e}ne \texttt{run\_distributed.sh} -- ta odpre \textbf{novo okno}, v katerem te\v{c}e porazdeljena simulacija.

\begin{vprasanje}
Zakaj se porazdeljeni na\v{c}in ne izvede kar v istem oknu?

\textbf{Odgovor:} Ker potrebuje \textbf{ve\v{c} lo\v{c}enih procesov}, ki jih mora zagnati \texttt{mpjrun.sh}. Obstoje\v{c}i proces se ne more sam ``razmno\v{z}iti'' v MPI skupino -- MPI okolje mora obstajati \v{z}e ob zagonu programa (\texttt{MPI.Init()} je prvi klic v \texttt{DistributedRunner}).
\end{vprasanje}

%======================================================================
\section{Pogosta vpra\v{s}anja za zagovor}
%======================================================================

\begin{vprasanje}
Kako zagotovite, da vse tri verzije dajo enake rezultate?

\textbf{Odgovor:} Vse tri verzije uporabljajo enak seed (42) za generiranje za\v{c}etnih delcev, isto ra\v{c}unsko jedro (\texttt{ForceKernel}) in enak dvofazni algoritem. Razlika je samo v tem, \textbf{kdo} izra\v{c}una katere sile -- rezultat je enak, ker je vsota sil neodvisna od vrstnega reda se\v{s}tevanja (do numeri\v{c}ne natan\v{c}nosti).
\end{vprasanje}

\begin{vprasanje}
Kaj je Amdahlov zakon in kako vpliva na va\v{s}o simulacijo?

\textbf{Odgovor:} Amdahlov zakon pravi, da je pohitritev omejena z dele\v{z}em kode, ki \textbf{mora} te\v{c}i sekvencno. V na\v{s}em primeru je Faza 2 (posodobitev, $O(n)$) sekvencna, a je zanemarljivo majhna v primerjavi s Fazo 1 ($O(n^2)$), zato se dobro skaliramo.
\end{vprasanje}

\begin{vprasanje}
Zakaj je porazdeljena verzija po\v{c}asnej\v{s}a od vzporedne?

\textbf{Odgovor:} Vzporedna verzija deli pomnilnik -- niti direktno berejo isti seznam delcev brez kopiranja. Porazdeljena verzija mora vsak cikel izvesti Allgatherv, ki kopira $6 \cdot n$ doublov med vsemi procesi. Ta komunikacijski overhead je \v{s}e posebej izrazit pri lokalnem zagonu (multicore), kjer ni dejanske mre\v{z}ne prednosti.
\end{vprasanje}

\begin{vprasanje}
Kdaj bi bila porazdeljena verzija hitrej\v{s}a od vzporedne?

\textbf{Odgovor:} Ko imamo \textbf{ve\v{c} ra\v{c}unalnikov} -- na primer 4 ra\v{c}unalnike z vsak 10 jedri. Vzporedna verzija je omejena na 10 niti enega ra\v{c}unalnika, porazdeljena pa bi uporabila $4 \times 10 = 40$ jeder.
\end{vprasanje}

\begin{vprasanje}
Kako bi izbolj\v{s}ali \v{c}asovno kompleksnost?

\textbf{Odgovor:} Z uporabo \textbf{prostorskega particioniranja} (Barnes-Hut algoritem s quadtree), ki zmanj\v{s}a kompleksnost iz $O(n^2)$ na $O(n \log n)$. Delce, ki so dale\v{c} stran, obravnavamo kot eno skupino namesto posamezno.
\end{vprasanje}

\begin{vprasanje}
Zakaj GUI te\v{c}e neodvisno od simulacije?

\textbf{Odgovor:} Izris poteka v lo\v{c}enem timerju (Swing Timer na EDT -- Event Dispatch Thread), simulacija pa v svoji niti. Tako GUI ne blokira izra\v{c}una in ostaja odziven. \v{C}e bi izra\v{c}un tekel na EDT, bi bil vmesnik zamrznjen.
\end{vprasanje}

\begin{vprasanje}
Zakaj se niti v ParallelSimulation ne ``stepejo na prste'' pri pisanju v forces[]?

\textbf{Odgovor:} Ker vsaka nit pi\v{s}e v \textbf{ekskluziven del} polja. Nit $t$ pi\v{s}e samo v indekse $[t \cdot \text{chunkSize}, (t+1) \cdot \text{chunkSize})$. Nobena druga nit ne pi\v{s}e v iste indekse. To je primer \textbf{prostorske dekompozicije} brez potrebe po zaklepanju.
\end{vprasanje}

\begin{vprasanje}
Kako zagotovite, da primerjate paralelizacijo in ne razlik med implementacijami?

\textbf{Odgovor:} Vsi trije na\v{c}ini kli\v{c}ejo \textbf{isti metodi} \texttt{ForceKernel.computeForces()} in \texttt{ForceKernel.integrate()} nad \textbf{isto podatkovno postavitvijo}, in te\v{c}ejo v \textbf{isti zanki} \texttt{SimulationRunner.runCycles()}. Razlikujejo se samo v mejah obmo\v{c}ja, ki ga dobi vsak izvajalec, in v na\v{c}inu sinhronizacije na koncu cikla. \v{C}e bi vsak na\v{c}in imel svojo kopijo fizike, bi izmerjena razmerja odra\v{z}ala razlike med temi kopijami in ne u\v{c}inka paralelizacije.
\end{vprasanje}

\begin{vprasanje}
Zakaj so rezultati vseh treh na\v{c}inov \textbf{popolnoma} enaki, ne le pribli\v{z}no?

\textbf{Odgovor:} Sile na delec $i$ se v vseh na\v{c}inih se\v{s}tevajo v \textbf{istem vrstnem redu} ($j = 0 \ldots n$), saj vsak izvajalec za svoj delec vedno pregleda vse delce po nara\v{s}\v{c}ajo\v{c}em indeksu. Ker se se\v{s}tevanje v plavajo\v{c}i vejici nikjer ne prerazporedi, ni razlik niti v zadnjem bitu. Razdelitev dela vpliva na to, \textbf{kdo} ra\v{c}una, ne na to, \textbf{v kak\v{s}nem vrstnem redu} se\v{s}teva.
\end{vprasanje}

\begin{vprasanje}
Zakaj pohitritev ne dose\v{z}e 10$\times$ na 10-jedrnem procesorju?

\textbf{Odgovor:} Apple M4 ima \textbf{4 zmogljiva (P) in 6 var\v{c}nih (E) jeder}. Prve \v{s}tiri niti dobijo P-jedra in prispevajo skoraj polno zmogljivost, naslednjih \v{s}est pa E-jedra, ki so bistveno po\v{c}asnej\v{s}a. Ker so kosi \textbf{enako veliki} in je na koncu vsakega cikla pregrada, vsi \v{c}akajo na najpo\v{c}asnej\v{s}ega. Zato u\v{c}inkovitost pri 4 izvajalcih (51--73\,\%) mo\v{c}no presega u\v{c}inkovitost pri 10 izvajalcih (22--29\,\%).
\end{vprasanje}

\end{document}
